
\begin{abstract}

  Modern GPU workloads in machine learning and HPC present complex behaviors that are challenging to observe, optimize and customize. An ideal GPU extension framework requires four key characteristics: (1) intra-kernel programmability to observe and modify running kernels at the instruction level, (2) safety to prevent crashes and memory corruption, (3) efficiency to minimize runtime overhead, and (4) unified programmability to coherently manage logic across the host CPU and GPU. We analyze existing approaches and find they fall short: host-side tools treat the GPU as a black box, failing on intra-kernel programmability; vendor-specific libraries like CUPTI and NVBit are GPU specific, lacking a unified programming model to control system-wide behavior and often introducing substantial overhead. This paper presents \sys, the first framework that extends the eBPF paradigm for general-purpose, programmable extensions directly into GPU kernels. \sys defines GPU-specific attach points for memory operations and synchronization primitives, verifies and JIT-compiles eBPF bytecode to PTX/SPIR-V with runtime injection, and provides on-device helper calls, plus shared CPU-GPU maps. By executing eBPF programs safely within GPU contexts, our framework enables a suite of bcc-style observability tools, programmable UVM management, and custom scheduling policies with 10× lower overhead than existing binary instrumentation. We demonstrate \sys's effectiveness through bcc-style GPU profiling tools (kernelretsnoop, threadhist, launchlate, mem\_trace) averaging 2-3\% overhead, UVM prefetching reducing page faults by 50-65\%, and fine-grain scheduling with sub-microsecond decision latency.
\end{abstract}


\section{Introduction}

The complexity of modern GPU architectures, with their SIMT execution model and deep memory hierarchies, makes understanding and optimizing workloads a critical challenge for developers~\cite{kousha2019designing,shen2018cudaadvisor,zhou2021measurement,stocks2024breaking,das2023large,fedeli2022pushing,liu2021closing}. For example, an ML engineer debugging tail latency in a model~\cite{guan2024amanda,kwon2023efficient} may be unable to distinguish between a CPU-side data bottleneck, a memory stall inside a specific GPU kernel, or inefficient scheduling between kernels. Moreover, a cloud provider may want to enforce fair resource scheduling inside a GPU, or dynamically offload computation to the CPU based on real-time power constraints without inference framework changes. An ideal solution requires a general-purpose extension framework that can safely inject custom logic for monitoring, optimization, and policy enforcement across the entire software stack.

Such a framework must meet four design goals. First, it requires intra-kernel programmability, the ability to observe and modify behavior at the warp or instruction level inside a running GPU kernel. Second, it requires safety, assuring that a buggy extension cannot crash the GPU or corrupt application memory. Third, it needs efficiency, executing at near-native speed to be viable in production. Finally, it demands unified programmability, allowing developers to implement and coordinate custom logic across the CPU and GPU within a single, coherent programming model, such as linking events across stacks.

Alas, existing frameworks do not meet these needs. Host-side observability tools, which attach probes to CUDA libraries or drivers, are not comprehensive; they treat the GPU as a black box, offering no way to inject custom logic into its execution~\cite{hong2017gpu,kato2011timegraph}. Conversely, vendor-specific tools like NVIDIA Nsight and CUPTI~\cite{cupti} provide deep visibility but are siloed; they do not integrate with general-purpose system frameworks like eBPF~\cite{gregg16,gregg2021computing,Bachl21,jia2023practical,Zhong22,yang2023lambda,jia2023programmable,ghigoff2021bmc,chen2025etran,xu2025state,benson2024netedit,zhou2024dint,wu2025vep,bpftime,mao2024merlin,jin2024beebox,wang2024seak,zhong2022xrp}, making it difficult to implement coordinated policies that span both host and device. Binary instrumentation tools like NVBit~\cite{villa2019nvbit} can inject logic but work only on vendor-specific GPU architectures and can impose substantial overhead for fine-grained instruction probes, often trading off safety and efficiency, and they remain vendor-specific, failing to provide a unified extension model.

This paper presents \sys, a new system that bridges this gap by extending the proven eBPF paradigm~\cite{eBPFWindows} for safe, dynamic, and efficient general-purpose extensions directly into the GPU context. \sys is based on two design principles. First, it maintains compatibility with the eBPF ecosystem, allowing developers to use familiar tools and programming models. Second, it resolves the programmability and performance gap by executing eBPF extensions directly on the device, without kernel or host involvement in the hot path. Realizing these principles requires overcoming three challenges: (1) defining meaningful attach points in a SIMT execution model, (2) enforcing eBPF's safety guarantees via on-device verification and sandboxing, and (3) implementing efficient and coherent eBPF maps across the CPU-GPU memory domain.

We implement \sys on Linux without requiring any Linux kernel modifications. Our framework introduces four key technical contributions: (1) GPU-specific eBPF attach points that define where programs can safely hook into kernel execution, memory operations~\cite{pankratz2021vulkan,sellers2016vulkan}, and synchronization primitives; (2) an eBPF-to-PTX JIT compiler that translates verified eBPF bytecode into NVIDIA's PTX representation for runtime injection; (3) CPU-GPU shared eBPF maps that enable efficient cross-domain state management; and (4) a PTX-level sandboxing mechanism that enforces eBPF's safety guarantees within the GPU's execution environment. Through these mechanisms, \sys{} allows developers to dynamically instrument, observe, and control GPU kernels without application source changes or GPU kernel interruption.

We demonstrate \sys's generality on use-cases that go beyond simple observation, including fine-grained memory tracing, programmable fair-share scheduling, and adaptive memory management for ML inference. These applications showcase how \sys enables novel GPU customization and optimization with significantly lower overhead than state-of-the-art alternatives.% The code is open-sourced under the \url{https://github.com/eunomia-bpf/eGPU} repo.

In sum, our contributions are:
\begin{itemize}[leftmargin=*,itemsep=0pt]
    \item An examination of GPU extension use-cases and the distillation of four design goals: intra-kernel programmability, safety, efficiency, and unified programmability.

    \item The design and implementation of \sys, a framework that brings the eBPF programming model into GPU kernels for general-purpose extensions via on-device JIT compilation, attach points, and runtime support.

    \item An evaluation showing that \sys enables novel GPU customization, optimization, and observability with significantly lower overhead (10× reduction) than state-of-the-art alternatives.

% \begin{itemize}[leftmargin=*,itemsep=0pt]
%   \item \textbf{MemTrace:} Per‐warp memory‐access tracing to detect cache‐miss patterns and bank conflicts inside active kernels, enabling pinpoint diagnosis of memory‐bound bottlenecks.
%   \item \textbf{Fair share Scheduling:} Real‐time collection of sharing information during execution to rebalance workloads and reduce tail latencies in irregular or unbalanced kernels.
%   \item \textbf{CPU–GPU Collaborative Inference (ML Green):} Sub‐100 µs end‐to‐end latency breakdowns across host copy routines and GPU kernels for large‐model pipelines (e.g., Llama 2), driving dynamic batch‐size adaptation and prefetch ordering.
% \end{itemize}
\end{itemize}


